\documentclass[12pt]{article}
\usepackage[utf8]{inputenc}
\usepackage[T1]{fontenc}
\usepackage{hyperref}
\title{Tagin Language Resources for K-12 Curriculum Development}
\author{Compiled by Claude Sonnet 5 (claude-sonnet-5) for Andr\'as Kornai}
\date{}
\begin{document}
\maketitle

\section{Introduction}
Tagin (ISO 639-3 tgj) is a Western Tani language of Sino-Tibetan, spoken in Arunachal Pradesh, India. It was historically oral and has recently been written in a Modified Roman Script. UNESCO lists it as Vulnerable. One news report says Tagin began to be taught in government schools in Tagin areas of Upper Subansiri. That report is undated in the search results, so the timing is unconfirmed.

\section{Primary Educational Resources}
No Tagin primers or textbooks were found as online resources. The Arunachal Times item below reports on school teaching: \url{https://arunachaltimes.in/?p=194393}

\section{Linguistic and Reference Materials}
\begin{itemize}
\item Tagin Living Dictionary (AI4Arunachal), with English and Hindi glosses, some audio, and a grammar section: \url{https://livingdictionaries.app/tagin}
\item About page for the living dictionary: \url{https://livingdictionaries.app/tagin/about}
\item UNESCO Atlas of the World's Languages page for Tagin. It lists print references: Das Gupta (1983), \emph{An Outline on Tagin Language}; Kepor Mara (2012), \emph{Illustrated Tagin-English Dictionary}; G.B. Baram (2013), \emph{Tagin Language: Grammar and Dictionary}; and a 2021 sociolinguistic survey by Abraham and Sako. \url{https://ar.wal.unesco.org/index.php/en/countries/india/languages/tagin}
\item Grambank source entry for Das Gupta's Tagin outline: \url{https://grambank.clld.org/sources/sDasGuptaTagin}
\item Wikipedia on the Bangni-Tagin language: \url{https://en.wikipedia.org/wiki/Bangni-Tagin_language}
\end{itemize}

\section{Religious and Community Resources}
None found.

\section{Digital and Audio Resources}
The living dictionary above includes some audio recordings.

\section{Curriculum Development Context}
The language is oral-first, with a newer Roman-based script and a community-built online dictionary. Naming is confusing: sources also use West Dafla and Bangni-Tagin, and the separate language Na has a different ISO code. Speaker counts differ across sources, with the 2011 census figure much higher than the figure given in the prompt.

\end{document}